\MathBookSourcePage{153}
\subsection{一般二阶椭圆算子}
\label{sec:ch05-general-second-order-elliptic-operators}
\setcounter{thmcount}{0}
\setcounter{equation}{0}

现在考虑更一般的椭圆算子及其相应变分形式. 设 $a_{ij}\in L^\infty(\Omega)$, $i,j=1,\ldots,n$, 且 $a_{ij}=a_{ji}$. 若存在正常数 $\alpha$, 使
\MBSourceItem{1}
\begin{equation}
\MathBookFormulaID{formula-ch05-uniform-ellipticity}
\label{eq:ch05-uniform-ellipticity}
\sum_{i,j=1}^na_{ij}(x)\xi_i\xi_j
\geq\alpha\sum_{i=1}^n\xi_i^2
\quad\forall\xi\in\mathbb R^n,
\quad\text{在 }\Omega\text{ 中几乎处处},
\end{equation}
则称系数矩阵 $(a_{ij})$ 一致椭圆, 即矩阵 $(a_{ij})$ 一致且几乎处处正定. 与椭圆系数矩阵相联系的椭圆算子定义为
\MBSourceItem{2}
\begin{equation}
\MathBookFormulaID{formula-ch05-symmetric-elliptic-operator}
\label{eq:ch05-symmetric-elliptic-operator}
Au(x):=-\sum_{i,j=1}^n\frac{\partial}{\partial x_j}
\left(a_{ij}(x)\frac{\partial u}{\partial x_i}(x)\right).
\end{equation}
例如, 若对所有 $x\in\Omega$, $(a_{ij}(x))$ 是单位矩阵, 就得到前面研究的 Laplace 算子, 且可取 $\alpha=1$. 与式~\eqref{eq:ch05-symmetric-elliptic-operator} 相联系的自然双线性形式为
\MBSourceItem{3}
\begin{equation}
\MathBookFormulaID{formula-ch05-symmetric-elliptic-bilinear-form}
\label{eq:ch05-symmetric-elliptic-bilinear-form}
a(u,v):=\int_\Omega\sum_{i,j=1}^na_{ij}(x)
\frac{\partial u}{\partial x_i}(x)
\frac{\partial v}{\partial x_j}(x)\,dx.
\end{equation}
注意并未假设系数具有任何光滑性, 因而允许不连续系数. 这不会给变分表述带来困难, 但式~\eqref{eq:ch05-symmetric-elliptic-operator} 中对不连续函数作微分需要进一步论证, 例如可通过变分表述本身. 在式~\eqref{eq:ch05-uniform-ellipticity} 中以 $\partial v/\partial x_i$ 代替 $\xi_i$ 并积分, 立即得到下述估计.

\MBSourceItem{4}
\begin{Lemma}
\label{lem:ch05-elliptic-energy-lower-bound}
设系数矩阵满足式~\eqref{eq:ch05-uniform-ellipticity}. 则相应双线性形式~\eqref{eq:ch05-symmetric-elliptic-bilinear-form} 满足
\[
\MathBookFormulaID{formula-ch05-elliptic-energy-lower-bound}
a(v,v)\geq\alpha|v|_{H^1(\Omega)}^2
\quad\forall v\in H^1(\Omega),
\]
其中 $\alpha$ 是式~\eqref{eq:ch05-uniform-ellipticity} 中的常数.
\end{Lemma}

与一般形式~\eqref{eq:ch05-symmetric-elliptic-bilinear-form} 相联系的自然边界条件比式~\eqref{eq:ch05-pure-neumann-condition} 更复杂. 在命题~\ref{prop:ch05-sobolev-integration-by-parts} 中取 $v:=-a_{ji}\partial u/\partial x_j$, 再对 $i$ 求和(交换 $i,j$ 后), 得
\[
\MathBookFormulaID{formula-ch05-general-green-identity}
\int_\Omega Auw\,dx
=a(u,w)-\int_{\partial\Omega}\sum_{i,j=1}^n
a_{ij}\frac{\partial u}{\partial x_j}\nu_iw\,ds.
\]
因此以算子形式表示的自然边界条件为
\MathBookSourcePage{154}
\MBSourceItem{5}
\begin{equation}
\MathBookFormulaID{formula-ch05-general-natural-boundary-condition}
\label{eq:ch05-general-natural-boundary-condition}
\sum_{i,j=1}^na_{ij}\frac{\partial u}{\partial x_j}\nu_i=0.
\end{equation}

变分形式~\eqref{eq:ch05-symmetric-elliptic-bilinear-form} 是对称的. 实际问题中一般也会出现非对称问题. 微分形式为
\MBSourceItem{6}
\begin{equation}
\MathBookFormulaID{formula-ch05-general-nonsymmetric-operator}
\label{eq:ch05-general-nonsymmetric-operator}
Au(x):=-\sum_{i,j=1}^n\frac{\partial}{\partial x_j}
\left(a_{ij}(x)\frac{\partial u}{\partial x_i}(x)\right)
+\sum_{k=1}^nb_k(x)\frac{\partial u}{\partial x_k}(x)+b_0(x)u(x)
\end{equation}
的问题可自然地用变分形式
\MBSourceItem{7}
\begin{equation}
\MathBookFormulaID{formula-ch05-general-nonsymmetric-form}
\label{eq:ch05-general-nonsymmetric-form}
a(u,v):=\int_\Omega\left(
\sum_{i,j=1}^na_{ij}\frac{\partial u}{\partial x_i}\frac{\partial v}{\partial x_j}
+\sum_{k=1}^nb_k\frac{\partial u}{\partial x_k}v+b_0uv
\right)dx
\end{equation}
表示. 下述定理部分回答了这种变分形式的强制性问题.

\MBSourceItem{8}
\begin{Theorem}[Gårding 不等式]
\label{thm:ch05-garding-inequality}
假设式~\eqref{eq:ch05-uniform-ellipticity} 成立, 且 $b_k\in L^\infty(\Omega)$, $k=0,\ldots,n$. 则存在常数 $K<\infty$, 使
\[
\MathBookFormulaID{formula-ch05-garding-inequality}
a(v,v)+K\|v\|_{L^2(\Omega)}^2
\geq\frac\alpha2\|v\|_{H^1(\Omega)}^2
\quad\forall v\in H^1(\Omega),
\]
其中 $\alpha$ 是式~\eqref{eq:ch05-uniform-ellipticity} 中的常数.
\end{Theorem}
这意味着, 虽然 $a(\cdot,\cdot)$ 本身未必具有强制性, 但加上充分大的常数倍 $L^2$ 内积后, 它在整个 $H^1(\Omega)$ 上具有强制性.
\begin{Proof}
由引理~\ref{lem:ch05-elliptic-energy-lower-bound},
\[
\MathBookFormulaID{formula-ch05-garding-proof-start}
a(v,v)+K\|v\|_{L^2(\Omega)}^2
\geq\alpha|v|_{H^1(\Omega)}^2
+\int_\Omega\left(\sum_{k=1}^nb_k(x)\frac{\partial v}{\partial x_k}(x)v(x)
+(b_0(x)+K)v(x)^2\right)dx.
\]
由 Hölder 不等式,
\[
\MathBookFormulaID{formula-ch05-first-order-term-bound}
\begin{aligned}
\left|\int_\Omega\sum_{k=1}^nb_k(x)\frac{\partial v}{\partial x_k}(x)v(x)\,dx\right|
&\leq\int_\Omega\sum_{k=1}^n|b_k(x)|
\left|\frac{\partial v}{\partial x_k}(x)\right||v(x)|\,dx\\
&\leq\sum_{k=1}^n\|b_k\|_{L^\infty(\Omega)}
\left\|\frac{\partial v}{\partial x_k}\right\|_{L^2(\Omega)}\|v\|_{L^2(\Omega)}\\
&\leq B|v|_{H^1(\Omega)}\|v\|_{L^2(\Omega)},
\end{aligned}
\]
其中
\MathBookSourcePage{155}
\MBSourceItem{9}
\begin{equation}
\MathBookFormulaID{formula-ch05-first-order-coefficient-bound}
\label{eq:ch05-first-order-coefficient-bound}
B^2:=\sum_{k=1}^n\|b_k\|_{L^\infty(\Omega)}^2.
\end{equation}
因此
\[
\MathBookFormulaID{formula-ch05-garding-proof-middle}
\begin{aligned}
a(v,v)+K\|v\|_{L^2(\Omega)}^2
&\geq\alpha|v|_{H^1(\Omega)}^2
-B|v|_{H^1(\Omega)}\|v\|_{L^2(\Omega)}
+\int_\Omega(b_0(x)+K)v(x)^2\,dx\\
&\geq\alpha|v|_{H^1(\Omega)}^2
-B|v|_{H^1(\Omega)}\|v\|_{L^2(\Omega)}
+(\beta+K)\|v\|_{L^2(\Omega)}^2,
\end{aligned}
\]
其中
\MBSourceItem{10}
\begin{equation}
\MathBookFormulaID{formula-ch05-essential-infimum-b0}
\label{eq:ch05-essential-infimum-b0}
\beta:=\operatorname*{ess\,inf}\{b_0(x):x\in\Omega\}.
\end{equation}
由算术--几何平均不等式~\eqref{eq:ch00-young-inequality-delta}, 得
\[
\MathBookFormulaID{formula-ch05-garding-proof-final}
a(v,v)+K\|v\|_{L^2(\Omega)}^2
\geq\frac\alpha2\left(|v|_{H^1(\Omega)}^2+\|v\|_{L^2(\Omega)}^2\right),
\]
只要
\MBSourceItem{11}
\begin{equation}
\MathBookFormulaID{formula-ch05-garding-k-condition}
\label{eq:ch05-garding-k-condition}
K\geq\frac\alpha2+\frac{B^2}{2\alpha}-\beta.
\end{equation}
注意若 $\beta>0$, $K$ 未必为正.
\end{Proof}

Gårding 不等式可解释如下. 给定椭圆算子 $A$, 存在常数 $K_0$, 使对所有 $K>K_0$, $A+K$ 的变分问题在整个 $H^1$ 上具有强制性, 从而式~\eqref{eq:ch05-mixed-boundary-conditions} 形式的任意边界条件都给出适定问题. 一般而言, 当 $K\leq K_0$ 时可能出现进一步约束, 如第~\ref{sec:ch05-pure-neumann-problem} 节研究的纯 Neumann 问题. 特别地, 可证明(Agmon 1965)存在一列没有有限聚点的值 $K_1>K_2>\cdots$(即 $A$ 的特征值), 使得只要对任意 $i$ 有 $K\ne K_i$, 与 $A+K$ 相联系的变分问题就有唯一解. $A+K_i$ 问题也可在有限个线性约束下求解, 这些约束与第~\ref{sec:ch05-pure-neumann-problem} 节中的约束相似. 此外, 若所有系数和边界足够光滑, 当边界条件为纯 Neumann 条件~\eqref{eq:ch05-pure-neumann-condition} 或纯 Dirichlet 条件(式~\eqref{eq:ch05-mixed-boundary-conditions} 中 $\Gamma=\partial\Omega$)时, 解光滑且
\MBSourceItem{12}
\begin{equation}
\MathBookFormulaID{formula-ch05-general-elliptic-regularity}
\label{eq:ch05-general-elliptic-regularity}
\|u\|_{W_p^k(\Omega)}
\leq C_{K,k,\Omega,A,p}\|(A+K)u\|_{W_p^{k-2}(\Omega)},
\qquad1<p<\infty.
\end{equation}
一般而言, 指数 $k$ 可能依赖于问题各组成部分的光滑程度; 若它们全为 $C^\infty$, 则 $k$ 可任意大. 当 $\Omega$ 是 $\mathbb R^2$ 中的凸多边形时, 已知式~\eqref{eq:ch05-general-elliptic-regularity} 对 $k\leq2$ 成立(Grisvard 1985).
